
\documentclass{article}
\usepackage{colortbl}
\usepackage{makecell}
\usepackage{multirow}
\usepackage{supertabular}

\begin{document}

\newcounter{utterance}

\centering \large Interaction Transcript for game `function\_detective', experiment `two\_numbers\_passive\_examples\_oneshot', episode 298 with gpt{-}5.4{-}mini.
\vspace{24pt}

{ \footnotesize  \setcounter{utterance}{1}
\setlength{\tabcolsep}{0pt}
\begin{supertabular}{c@{$\;$}|p{.15\linewidth}@{}p{.15\linewidth}p{.15\linewidth}p{.15\linewidth}p{.15\linewidth}p{.15\linewidth}}
    \# & \multicolumn{2}{c}{Player} && \multicolumn{2}{c}{Game Master} \\
    \hline

    \theutterance \stepcounter{utterance}  
    & & & \multicolumn{4}{p{0.6\linewidth}}{
        \cellcolor[rgb]{0.9,0.9,0.9}{
            \makecell[{{p{\linewidth}}}]{
                \texttt{\tiny{[P1$\langle$GM]}}
                \texttt{I have a hidden Python function for you to figure out.} \\
\\ 
\texttt{Your job is to infer what it does and then submit a correct Python implementation.} \\
\\ 
\texttt{Function signature:} \\
\texttt{(a: int, b: int) {-}> int} \\
\\ 
\texttt{You have 1 turns total.} \\
\\ 
\texttt{I have already given you a fixed list of valid input{-}output examples. Solve using those examples only.} \\
\\ 
\texttt{You must make your best SOLVE attempt no later than the second{-}to{-}last turn.} \\
\texttt{Do not wait until the very end.} \\
\\ 
\texttt{You can use these actions:} \\
\texttt{{-} JSON format described at the end of the prompt} \\
\\ 
\\ 
\texttt{Do not add import statements in your guess.} \\
\texttt{Anything you might need is already available, so just write the function body in hypothesis and return the answer.} \\
\\ 
\texttt{Preloaded examples:} \\
\texttt{1. ((1, 10), 3)} \\
\texttt{2. (({-}10, 1), 4)} \\
\texttt{3. ((10, {-}99), 4)} \\
\texttt{4. (({-}99, 0), 0)} \\
\texttt{5. (({-}1000, 0), 0)} \\
\texttt{6. (({-}2, {-}1), 2)} \\
\texttt{7. ((99, 0), 0)} \\
\texttt{8. (({-}1000, {-}99), 6)} \\
\texttt{9. (({-}99, 2), 5)} \\
\texttt{10. ((1000, {-}2), 2)} \\
\\ 
\texttt{Scoring:} \\
\texttt{{-} correctness is required for any non{-}zero score} \\
\texttt{{-} in interactive modes, solving with fewer information requests is better} \\
\\ 
\texttt{Input formatting reminders:} \\
\texttt{{-} strings must be quoted} \\
\texttt{{-} booleans must be True/False} \\
\texttt{{-} use commas between multiple values} \\
\texttt{{-} do not write any extra action tags in the same message} \\
\\ 
\texttt{Format example:} \\
\texttt{You write:} \\
\texttt{\{'current\_mode': <probing/confirming/don't know; choose whichever best describes your current state in the task>, 'input\_rationale': <brief description of why you chose this hypothesis>, 'hypothesis': <either the exact string "unknown" if you have no guess yet OR your current best guess as runnable Python matching the given signature, with a one{-}line docstring stating what it does. Never write placeholder code, if you cannot commit to a guess, use "unknown". Match the given signature.} \\
\texttt{e.g.} \\
\texttt{def f(x):} \\
\texttt{	"""Returns the absolute value of x."""} \\
\texttt{	return x if x >= 0 else {-}x} \\
\\ 
\texttt{Use "unknown" if you have no guess.>, 'confidence\_score': <score between 1{-}100. 100 means highest confidence in your hypothesis>\}} \\
\\ 
\texttt{Output only one valid json per turn.} \\
            }
        }
    }
     \\ \\

    \theutterance \stepcounter{utterance}  
    & \multicolumn{4}{p{0.6\linewidth}}{
        \cellcolor[rgb]{0.561,0.737,0.561}{
            \makecell[{{p{\linewidth}}}]{
                \texttt{\tiny{[P1$\rangle$GM]}}
                \texttt{\{"current\_mode":"probing","input\_rationale":"The examples look mostly like a digit{-}counting function with special handling for sign/zeros: zero in b returns 0, same{-}sign pairs behave like total decimal digit counts of absolute values, and mixed{-}sign cases appear to add or remove based on the negative operand. I’m using a best{-}fit heuristic from the examples.","hypothesis":"def f(a: int, b: int) {-}> int:\textbackslash{}n    \textbackslash{}"\textbackslash{}"\textbackslash{}"Returns a sign{-}aware digit{-}count based on the two integers.\textbackslash{}"\textbackslash{}"\textbackslash{}"\textbackslash{}n    if b == 0:\textbackslash{}n        return 0\textbackslash{}n\textbackslash{}n    def digits(n: int) {-}> int:\textbackslash{}n        return len(str(n))\textbackslash{}n\textbackslash{}n    def nonzero\_digits(n: int) {-}> int:\textbackslash{}n        s = str(n)\textbackslash{}n        c = 0\textbackslash{}n        for ch in s:\textbackslash{}n            if ch != '0':\textbackslash{}n                c += 1\textbackslash{}n        return c\textbackslash{}n\textbackslash{}n    def trailing\_zeros(n: int) {-}> int:\textbackslash{}n        s = str(n)\textbackslash{}n        i = len(s) {-} 1\textbackslash{}n        c = 0\textbackslash{}n        while i >= 0 and s{[}i{]} == '0':\textbackslash{}n            c += 1\textbackslash{}n            i {-}= 1\textbackslash{}n        return c\textbackslash{}n\textbackslash{}n    aa = abs(a)\textbackslash{}n    bb = abs(b)\textbackslash{}n    base = digits(aa) + digits(bb)\textbackslash{}n\textbackslash{}n    if a < 0 < b:\textbackslash{}n        return base + nonzero\_digits(aa)\textbackslash{}n    if a > 0 > b and bb < 10:\textbackslash{}n        return base {-} trailing\_zeros(aa)\textbackslash{}n    return base","confidence\_score":22\}} \\
            }
        }
    }
    & & \\ \\

    \theutterance \stepcounter{utterance}  
    & & & \multicolumn{2}{p{0.3\linewidth}}{
        \cellcolor[rgb]{0.9,0.9,0.9}{
            \makecell[{{p{\linewidth}}}]{
                \texttt{\tiny{[GM$|$GM]}}
                \texttt{Game Verdict: LOSS} \\
            }
        }
    }
    & & \\ \\

    \theutterance \stepcounter{utterance}  
    & & & \multicolumn{2}{p{0.3\linewidth}}{
        \cellcolor[rgb]{0.9,0.9,0.9}{
            \makecell[{{p{\linewidth}}}]{
                \texttt{\tiny{[GM$|$GM]}}
                \texttt{Incorrect.} \\
\texttt{Failed Cases:} \\
\texttt{Input: {[}1000, 10{]} | Expected: 4 | Got: 6} \\
\texttt{Input: {[}2, {-}1{]} | Expected: 5 | Got: 2} \\
\texttt{Input: {[}1, {-}10{]} | Expected: 4 | Got: 3} \\
            }
        }
    }
    & & \\ \\

    \theutterance \stepcounter{utterance}  
    & & & \multicolumn{2}{p{0.3\linewidth}}{
        \cellcolor[rgb]{0.9,0.9,0.9}{
            \makecell[{{p{\linewidth}}}]{
                \texttt{\tiny{[GM$|$GM]}}
                \texttt{def product\_mod\_seven(a: int, b: int) {-}> int:} \\
\texttt{    return (a * b) \% 7} \\
            }
        }
    }
    & & \\ \\

    \theutterance \stepcounter{utterance}  
    & & & \multicolumn{2}{p{0.3\linewidth}}{
        \cellcolor[rgb]{0.9,0.9,0.9}{
            \makecell[{{p{\linewidth}}}]{
                \texttt{\tiny{[GM$|$GM]}}
                \texttt{def f(a: int, b: int) {-}> int:} \\
\texttt{    """Returns a sign{-}aware digit{-}count based on the two integers."""} \\
\texttt{    if b == 0:} \\
\texttt{        return 0} \\
\\ 
\texttt{    def digits(n: int) {-}> int:} \\
\texttt{        return len(str(n))} \\
\\ 
\texttt{    def nonzero\_digits(n: int) {-}> int:} \\
\texttt{        s = str(n)} \\
\texttt{        c = 0} \\
\texttt{        for ch in s:} \\
\texttt{            if ch != '0':} \\
\texttt{                c += 1} \\
\texttt{        return c} \\
\\ 
\texttt{    def trailing\_zeros(n: int) {-}> int:} \\
\texttt{        s = str(n)} \\
\texttt{        i = len(s) {-} 1} \\
\texttt{        c = 0} \\
\texttt{        while i >= 0 and s{[}i{]} == '0':} \\
\texttt{            c += 1} \\
\texttt{            i {-}= 1} \\
\texttt{        return c} \\
\\ 
\texttt{    aa = abs(a)} \\
\texttt{    bb = abs(b)} \\
\texttt{    base = digits(aa) + digits(bb)} \\
\\ 
\texttt{    if a < 0 < b:} \\
\texttt{        return base + nonzero\_digits(aa)} \\
\texttt{    if a > 0 > b and bb < 10:} \\
\texttt{        return base {-} trailing\_zeros(aa)} \\
\texttt{    return base} \\
            }
        }
    }
    & & \\ \\

    \theutterance \stepcounter{utterance}  
    & & & \multicolumn{2}{p{0.3\linewidth}}{
        \cellcolor[rgb]{0.9,0.9,0.9}{
            \makecell[{{p{\linewidth}}}]{
                \texttt{\tiny{[GM$|$GM]}}
                \texttt{Mode: passive\_examples\_oneshot} \\
\texttt{Accuracy: 0.310} \\
\texttt{Efficiency: 1.000} \\
\texttt{Parse errors: 0} \\
\texttt{Runtime errors: 0} \\
            }
        }
    }
    & & \\ \\

\end{supertabular}
}

\end{document}
